\documentclass{article}

\usepackage[margin=0.5in]{geometry}
\usepackage{amsmath, mathtools, array, booktabs, hyperref, amssymb}

\title{HW 6}

\author{Darin Peries}

\date{\today}

\begin{document}
\maketitle

\section*{Problem 1}

The two point boundary value problem: 

$$-\phi'' + (4x^2 + 2)\phi = 0, \ \phi(0)=1, \ \phi(1)=e$$

\noindent With solution $\phi(x) = e^{x^2}$.

\noindent Using the fourth order finite-difference scheme, the resulting matrix is:

\[
    A = \begin{pmatrix}
        30 + 12h^2(4x^2+2) & -16 & 1 & 0 & 0 & \dots \\
        -16 & 30 + 12h^2(4x^2+2) & -16 & 1 & 0 & 0 & \dots \\
        1 & -16 & 30 + 12h^2(4x^2+2) & -16 & 1 & 0 & 0 & \dots \\ 
        0 & 1 & -16 & 30 + 12h^2(4x^2+2) & -16 & 1 & 0 & 0 & \dots\\
        \vdots \ddots
    \end{pmatrix}
\]

\subsection*{Error Tables}

\begin{table}[h!]
\centering
\caption{Error results using the Thomas algorithm.}
\begin{tabular}{|c c c c c|}
\hline
$h$ &
$\|u_h-\phi_h\|_\infty$ &
$\dfrac{\|u_h-\phi_h\|_\infty}{h^3}$ &
$\dfrac{\|u_h-\phi_h\|_\infty}{h^4}$ &
$\dfrac{\|u_h-\phi_h\|_\infty}{h^5}$ \\
\hline
$0.5$ &
$6.8077066\times10^{-2}$ &
$5.4461653\times10^{-1}$ &
$1.0892331$ &
$2.1784661$ \\

$0.25$ &
$1.6671510\times10^{-2}$ &
$1.0669766$ &
$4.2679065$ &
$1.7071626\times10^{1}$ \\

$0.125$ &
$2.0444497\times10^{-3}$ &
$1.0467582$ &
$8.3740659$ &
$6.6992527\times10^{1}$ \\

$0.0625$ &
$1.8761033\times10^{-4}$ &
$7.6845190\times10^{-1}$ &
$1.2295230\times10^{1}$ &
$1.9672369\times10^{2}$ \\

$0.03125$ &
$1.4352246\times10^{-5}$ &
$4.7029440\times10^{-1}$ &
$1.5049421\times10^{1}$ &
$4.8158146\times10^{2}$ \\

$0.015625$ &
$1.0286328\times10^{-6}$ &
$2.6964993\times10^{-1}$ &
$1.7257595\times10^{1}$ &
$1.1044861\times10^{3}$ \\
\hline
\end{tabular}
\end{table}

\begin{table}[h!]
\centering
\caption{Error results using gradient descent.}
\begin{tabular}{|c c c c c|}
\hline
$h$ &
$\|u_h-\phi_h\|_\infty$ &
$\dfrac{\|u_h-\phi_h\|_\infty}{h^3}$ &
$\dfrac{\|u_h-\phi_h\|_\infty}{h^4}$ &
$\dfrac{\|u_h-\phi_h\|_\infty}{h^5}$ \\
\hline
$0.5$ &
$6.8077066\times10^{-2}$ &
$5.4461653\times10^{-1}$ &
$1.0892331$ &
$2.1784661$ \\

$0.25$ &
$1.6671510\times10^{-2}$ &
$1.0669766$ &
$4.2679065$ &
$1.7071626\times10^{1}$ \\

$0.125$ &
$2.0444496\times10^{-3}$ &
$1.0467582$ &
$8.3740656$ &
$6.6992525\times10^{1}$ \\

$0.0625$ &
$1.8761023\times10^{-4}$ &
$7.6845151\times10^{-1}$ &
$1.2295224\times10^{1}$ &
$1.9672359\times10^{2}$ \\

$0.03125$ &
$1.4352103\times10^{-5}$ &
$4.7028972\times10^{-1}$ &
$1.5049271\times10^{1}$ &
$4.8157667\times10^{2}$ \\

$0.015625$ &
$1.0282268\times10^{-6}$ &
$2.6954349\times10^{-1}$ &
$1.7250784\times10^{1}$ &
$1.1040502\times10^{3}$ \\
\hline
\end{tabular}
\end{table}

\newpage

\begin{table}[h!]
\centering
\caption{Error results using the conjugate gradient method.}
\begin{tabular}{|c c c c c|}
\hline
$h$ &
$\|u_h-\phi_h\|_{\infty}$ &
$\dfrac{\|u_h-\phi_h\|_{\infty}}{h^3}$ &
$\dfrac{\|u_h-\phi_h\|_{\infty}}{h^4}$ &
$\dfrac{\|u_h-\phi_h\|_{\infty}}{h^5}$ \\
\hline
$0.5$ &
$6.8077066 \times 10^{-2}$ &
$5.4461653 \times 10^{-1}$ &
$1.0892331$ &
$2.1784661$ \\

$0.25$ &
$1.6671510 \times 10^{-2}$ &
$1.0669766$ &
$4.2679065$ &
$1.7071626 \times 10^{1}$ \\

$0.125$ &
$2.0444497 \times 10^{-3}$ &
$1.0467582$ &
$8.3740658$ &
$6.6992527 \times 10^{1}$ \\

$0.0625$ &
$1.8761032 \times 10^{-4}$ &
$7.6845186 \times 10^{-1}$ &
$1.2295230 \times 10^{1}$ &
$1.9672368 \times 10^{2}$ \\

$0.03125$ &
$1.4352247 \times 10^{-5}$ &
$4.7029444 \times 10^{-1}$ &
$1.5049422 \times 10^{1}$ &
$4.8158151 \times 10^{2}$ \\

$0.015625$ &
$1.0286984 \times 10^{-6}$ &
$2.6966712 \times 10^{-1}$ &
$1.7258696 \times 10^{1}$ &
$1.1045565 \times 10^{3}$ \\
\hline
\end{tabular}
\end{table}

\subsection*{Error Plots}

\begin{figure}[h!]
    \centering
    \includegraphics[width=0.8\linewidth]{thomas_errors.png}
    \caption{Errors from the Thomas Algorithm}
\end{figure}

\begin{figure}[h!]
    \centering
    \includegraphics[width=0.8\linewidth]{cg_errors.png}
    \caption{Errors from the Conjugate Gradient Algorithm}
\end{figure}

\begin{figure}[h!]
    \centering
    \includegraphics[width=0.8\linewidth]{gd_errors.png}
    \caption{Errors from the Gradient Descent Algorithm}
\end{figure}

\newpage

\subsection*{Discussion}

The main trend from the columns comes from the fact that as $h$ decreases, the error also decreases.
When dividing by the grid spacing, the error will converge to some constant based on the order of the discritization method.
Based on the results, the behavior trends towards an $h^4$ order of approximation.

\section*{Problem 2}

The temperature distribution $\Phi$ on a metal plate $[0,1]\times[0,1]$ with no heat source is given by the equation:

$$\Phi_{xx} + \Phi_{yy} = 0, \ (x, \ y)\in (0, \ 1) \times (0,\ 1)$$

\noindent The matrix created from the five-point stencil is:

$$A_h = I_m \otimes T_m + T_m \otimes I_m$$

\[
    T_m = \begin{pmatrix}
            2 & -1 & 0 & \dots & 0 & 0 \\
            -1 & 2 & -1 & 0 & \dots & 0\\
            0 & -1 & 2 & -1 & 0 & \dots & \\
            \vdots & \ddots
        \end{pmatrix}
\]

\subsection*{Residual Tables}

\begin{table}[h!]
\centering
\caption{Jacobi iteration results for the two-dimensional Laplace problem.}
\begin{tabular}{|c c c c|}
\hline
$h$ & $k$ &
$\displaystyle \frac{\|r_h^{(k)}\|_\infty}
{\|r_h^{(k-1)}\|_\infty}$ &
$\rho(B_J)$ \\
\hline
$1/4$  & 11 & 0.500488 & 0.707107 \\
$1/8$  & 32 & 0.971291 & 0.923880 \\
$1/16$ & 65 & 0.958709 & 0.980785 \\
\hline
\end{tabular}
\end{table}

\begin{table}[h!]
\centering
\caption{Gauss--Seidel iteration results for the two-dimensional Laplace problem.}
\begin{tabular}{|c c c c|}
\hline
$h$ & $k$ &
$\displaystyle \frac{\|r_h^{(k)}\|_\infty}
{\|r_h^{(k-1)}\|_\infty}$ &
$\rho(B_{GS})$ \\
\hline
$1/4$  & 7  & 0.502953 & 0.500000 \\
$1/8$  & 19 & 0.857683 & 0.853553 \\
$1/16$ & 36 & 0.955950 & 0.961940 \\
\hline
\end{tabular}
\end{table}

\begin{table}[h!]
\centering
\caption{Optimal SOR iteration results for the two-dimensional Laplace problem.}
\begin{tabular}{|c c c c c|}
\hline
$h$ & $\omega_\ast$ & $k$ &
$\displaystyle \frac{\|r_h^{(k)}\|_\infty}
{\|r_h^{(k-1)}\|_\infty}$ &
$\rho(B_{\omega_\ast})$ \\
\hline
$1/4$  & 1.171573 & 5  & 0.068282 & 0.171573 \\
$1/8$  & 1.446463 & 11 & 0.649997 & 0.446463 \\
$1/16$ & 1.673514 & 23 & 0.840870 & 0.673514 \\
\hline
\end{tabular}
\end{table}

\newpage

\subsection*{Consecutive residual Plot}

\begin{figure}[h!]
    \centering
    \includegraphics[width=0.8\linewidth]{residual_ratios.png}
    \caption{Dashed lines are $\rho(B)$}
\end{figure}

\subsection*{Discussion}

As $h$ increases, the final consecutive residual ratio decreases, which is seen across all algorithms with the slight exception of the 
Jacobi method. The consecutive ratio for the Gauss-Seidel method is almost identical to the spectral radius of $B_{GS}$, since the spectral radius is much smaller
and the method converges at a faster rate. The Jacobi method has a larger spectral radius, so the ratio is greater than or equal to the Gauss-Seidel.
Lastly, the SOR has the smallest spectral radius and the consecutive residual ratio is the smallest.

\newpage

\subsection*{Iterative Methods Consecutive-Residual plots}

\begin{figure}[h!]
    \centering
    \includegraphics[width=0.8\linewidth]{gd_cg_residual_ratios.png}
\end{figure}

\noindent The initial guess was $x^{(0)}=0$ and the tolerance is $\epsilon = 1e-8$.

\end{document}
